\section{Related Work}
\label{sec:related}
\paragraph{Task-aware visual communication.}
Learned image compression jointly models a latent representation and its
coding distribution; scale hyperpriors are a representative
approach~\cite{hyperprior}. LimitNet studies progressive content-aware
offloading for constrained devices and networks~\cite{limitnet}.
For VLA policies, SPARC trains an image reconstruction path using action
distortion, temporal masking, and a tilted rate penalty~\cite{sparc}.
Our work shares the task-oriented objective but recovers policy features
rather than images. Action awareness itself is not the distinguishing
claim; the communication interface and its division of computation are.

\paragraph{Collaborative inference and feature compression.}
Neurosurgeon studies partitioning inference between a mobile device and
the cloud~\cite{neurosurgeon}. BottleNet++ explicitly compresses
intermediate features for device--edge inference~\cite{bottlenetpp}.
RoboECC and RAPID extend collaborative deployment and partitioned
inference to VLA models~\cite{roboecc,rapid}. Our contribution is not
model splitting or activation transmission alone. We combine a compact policy-supervised visual path
with quantized temporal state and a concrete entropy-coded interface;
a controlled comparison with simple feature quantization remains part
of the evaluation needed to attribute its gains.

\paragraph{VLA interfaces and timely execution.}
OpenVLA and $\pi_{0.5}$ provide the policy families underlying our
implementation~\cite{openvla,pi05}. OpenVLA-OFT studies adaptation,
continuous action prediction, and action chunking~\cite{openvlaoft}.
Our method freezes the associated policy and changes its visual
input path. Distillation supplies the general teacher--student
principle~\cite{distillation}, while the per-policy visual interface
determines what representation the receiver must reconstruct.

ARIA studies inference across heterogeneous mobile
processors~\cite{aria}, and TimelyLLM coordinates generation with
physical agent execution~\cite{timelyllm}. Their settings reinforce
the need to relate compute timing to application behavior. The current
system uses existing asynchronous execution; its contribution concerns
visual representation and transmission rather than a new scheduling
algorithm.
